\section{Evidence Conditions and the Integrity/Sufficiency Axis}
\label{sec:conditions}

Each issue is evaluated under five conditions (Table~\ref{tab:conditions},
Figure~\ref{fig:concept}). ``T2'' is a static, single-shot surface; ``T3''
permits active tool calls. ``D'' denotes the tier-D (restricted) artifact and
``A'' the tier-A (decisive) artifact.

\begin{table}[t]
\centering
\small
\caption{The five evidence conditions per issue. Tier-D outcomes are the
post-patch F2P/P2P verdicts withheld on restricted surfaces; tier-A outcomes
are decisive.}
\label{tab:conditions}
\begin{tabularx}{\linewidth}{@{}l l l l X@{}}
\toprule
\textbf{Condition} & \textbf{Access} & \textbf{Tier-D outcomes} & \textbf{Tier-A outcomes} & \textbf{Condition gold} \\
\midrule
T2D-T            & static               & restricted (withheld) & absent              & \NEED \\
T3D-Strict       & active tools, strict & withheld even on fetch & absent             & \NEED \\
T3D-Escalate     & active, may escalate & granted after escalation & ---              & $=\gO$ \\
T2A-T            & static               & ---                    & tier-A sufficient  & $=\gO$ \\
T3A-T            & active               & ---                    & tier-A sufficient  & $=\gO$ \\
\bottomrule
\end{tabularx}
\end{table}

\begin{figure}[t]
\centering
\begin{tikzpicture}[x=1cm,y=1cm,
  verdict/.style={draw,rounded corners=2pt,line width=0.9pt,align=left,
                 inner sep=5pt,text width=4.55cm,minimum height=1.15cm},
  tier/.style={draw,rounded corners=2pt,line width=0.8pt,align=left,
                 inner sep=6pt,text width=6.55cm,minimum height=2.2cm},
  cond/.style={draw,rounded corners=2pt,line width=0.8pt,align=center,
                 inner sep=3pt,text width=2.55cm,minimum height=1.35cm},
  lbl/.style={font=\scriptsize\itshape},
  >={Stealth[length=2.4mm]}]

% --- three-way verdict ---
\node[verdict,draw=grAccept,text=grAccept,font=\small\bfseries]
  (acc) at (2.35,0) {ACCEPT\\[1pt]\lbl{\textcolor{grMuted}{patch acceptable under the granted scope}}};
\node[verdict,draw=grReject,text=grReject,font=\small\bfseries]
  (rej) at (7.85,0) {REJECT\\[1pt]\lbl{\textcolor{grMuted}{a decisive failure is evidenced}}};
\node[verdict,draw=grNeed,text=grNeed,font=\small\bfseries]
  (need) at (13.35,0) {NEED\_EVIDENCE\\[1pt]\lbl{\textcolor{grMuted}{available evidence is non-decisive}}};

% --- two gold tiers ---
\node[tier,fill=grOracle,draw=grNeed,anchor=north west] (og) at (0,-1.5)
  {\textcolor{grNeed}{\bfseries Oracle Gold}\\[2pt]\normalsize
   \textcolor{grInk}{\scriptsize ground truth from real pytest: F2P\,/\,P2P\,/\,symptom;
   ACCEPT if declared tests pass with no in-scope regression, REJECT on a decisive
   declared failure.}\\[3pt]
   \lbl{\textcolor{grMuted}{Scoped, not whole-project correctness.}}};
\node[tier,fill=grCond,draw=grReject,anchor=north west] (cg) at (7.15,-1.5)
  {\textcolor{grReject}{\bfseries Condition Gold}\\[2pt]\normalsize
   \textcolor{grInk}{\scriptsize the decision the granted evidence can justify:
   restricted surfaces $\rightarrow$ NEED\_EVIDENCE even when oracle is ACCEPT/REJECT;
   sufficient $\rightarrow$ oracle.}\\[3pt]
   \lbl{\textcolor{grMuted}{Distinguishes integrity from sufficiency.}}};

% --- condition ladder label ---
\node[anchor=north west,font=\small\bfseries,text=grInk] at (0,-5.15)
  {Evidence conditions (per issue)};

% --- five condition boxes ---
\node[cond,draw=grReject,anchor=north west] (c1) at (0,-5.75)
  {\textbf{T2D-T}\\[1pt]\scriptsize static restricted\\[2pt]\textcolor{grReject}{\itshape gold: NEED}};
\node[cond,draw=grReject,anchor=north west] (c2) at (3.16,-5.75)
  {\textbf{T3D-Strict}\\[1pt]\scriptsize active, outcomes withheld\\[2pt]\textcolor{grReject}{\itshape gold: NEED}};
\node[cond,draw=grAccept,anchor=north west] (c3) at (6.32,-5.75)
  {\textbf{T3D-Escalate}\\[1pt]\scriptsize active, then outcomes granted\\[2pt]\textcolor{grAccept}{\itshape gold: $=$ oracle}};
\node[cond,draw=grAccept,anchor=north west] (c4) at (9.48,-5.75)
  {\textbf{T2A-T}\\[1pt]\scriptsize static sufficient (tier-A)\\[2pt]\textcolor{grAccept}{\itshape gold: $=$ oracle}};
\node[cond,draw=grAccept,anchor=north west] (c5) at (12.64,-5.75)
  {\textbf{T3A-T}\\[1pt]\scriptsize active sufficient (tier-A)\\[2pt]\textcolor{grAccept}{\itshape gold: $=$ oracle}};

\draw[grMuted,line width=1pt,->] (c1) -- (c2);
\draw[grMuted,line width=1pt,->] (c2) -- (c3);
\draw[grMuted,line width=1pt,->] (c3) -- (c4);
\draw[grMuted,line width=1pt,->] (c4) -- (c5);
\end{tikzpicture}

\caption{Three-way review decision, the two gold tiers, and the five-condition
evidence ladder. Restricted conditions (red) warrant \NEED; sufficient
conditions (green) warrant the Oracle Gold decision.}
\label{fig:concept}
\end{figure}

T3D-\textbf{Strict} is the key insufficiency probe: an active agent can inspect
linkage, freshness, and coverage and read the symptom, but post-patch F2P/P2P
verdicts are structurally unreachable, so accepting or rejecting necessarily
upgrades non-decisive evidence. T3D-\textbf{Escalate} grants the decisive
outcomes only after the agent explicitly requests them. The evidence-serving
protocol is summarised in Algorithm~\ref{alg:protocol}.

\begin{algorithm}[t]
\caption{GroundedReview cell protocol (one issue $i$, condition $c$, model $m$)}
\label{alg:protocol}
\begin{algorithmic}[1]
\State \textbf{input}: issue text, candidate diff, condition $c$, served model id
\State open a session bound to the candidate patch identity \Comment{no gold field is ever served}
\State $d \gets{}$ decision policy for $c$ \Comment{strict, escalate, or static}
\If{$c$ is an active (T3) condition}
  \Loop
    \State model chooses: \textsc{getArtifact}, \textsc{verifyLinkage}, \textsc{getRunLog}, or \textsc{finish}
    \If{requested field is tier-D and policy is \textsc{strict}}
       \State reply \textbf{403 forbidden} \Comment{outcomes structurally unreachable}
    \ElsIf{policy is \textsc{escalate} without a stated reason}
       \State reply \textbf{400}; require a justification
    \ElsIf{escalation justified}
       \State reply \textbf{200} with header \texttt{X-Evidence-Tier-Escalation} and the decisive outcomes
    \EndIf
  \EndLoop
\EndIf
\State persist the \textbf{raw} response before any analysis \Comment{raw-first}
\State \textbf{output}: one of \ACCEPT, \REJECT, \NEED
\State score against condition gold $\gC(i,c)$; the issue, not the cell, is the unit
\end{algorithmic}
\end{algorithm}

\paragraph{Non-interference.}
Restricted projections are machine-scanned for post-patch verdicts, hidden
gold, acceptance digests, and reconstructing keys. We disclose one residual
instrument gap found in the corrective replication: a \emph{pre-patch}
baseline P2P map survived the tier-D projection. It cannot reconstruct gold
and was seen by none of the error cells, and it is mandated for removal before
the next live run (Section~\ref{sec:threats}). The corrected v2 projection used
for the prepared validation bed removes this field on every tool path
(Section~\ref{sec:validation}).
